\documentclass[portuguese]{technical}

\setdoctitle{[Relatório de Análise] A lista de projetos empresariais não é %
  exibida na página de criação do ECS na 8.5.1 no site do Brasil}
\setreportversion{HCS 8.5.1}
\setreportdate{2025-08-13}
\setdocversion{3.6.0}
\setreportscenario{Cenário Padrão}
\setheadertitle{Huawei Cloud --- Relatório Técnico}

\begin{document}
\makecover
\maketoc
\startbody

\begin{infobox}
Este documento é uma amostra do \textbf{modelo de Relatório Técnico Huawei} (\inlinecode{technical.cls}).
Ele demonstra a estrutura de 5 seções (problema, análise de causa raiz, causa raiz, condição de gatilho, solução alternativa)
junto com comandos de formatação compartilhados: caixas de destaque, blocos de código, tabelas, imagens e comandos inline.
\end{infobox}

\begin{problem}
No site do Brasil do HCS 8.5.1, o projeto empresarial padrão
não é exibido na página de criação do ECS. Isso impede que os
usuários selecionem o projeto empresarial correto ao provisionar
novas instâncias ECS.
\end{problem}

\begin{rootcauseanalysis}
\begin{enumerate}
  \item Analisadas as permissões do usuário que não visualiza o
    projeto empresarial padrão
  \item Verificado o comportamento do campo \inlinecode{auth\_action} na versão 8.5.1
  \item Confirmado que a função Server Administrator não possui
    permissões de ação refinadas
\end{enumerate}
\end{rootcauseanalysis}

\begin{rootcause}
Na versão 8.5.1, o cenário de permissão refinada foi otimizado.
Quando um ECS é provisionado, apenas os projetos empresariais com
as permissões apropriadas são listados. A função Server
Administrator não inclui ECS FullAccess.
\end{rootcause}

\begin{triggercondition}
Ocorre quando um usuário com apenas a função Server Administrator
tenta criar uma instância ECS no HCS 8.5.1.
\end{triggercondition}

\begin{workaround}

  \begin{impact}
  Para conceder ECS FullAccess a um usuário, é necessário modificar
  a permissão do grupo de usuários. Esta permissão é a permissão de
  administrador do ECS e está incluída na permissão do Server
  Administrator.
  \end{impact}

  \begin{backupdata}
  N/A --- nenhuma modificação de dados necessária.
  \end{backupdata}

  \begin{workaroundsteps}
  \begin{enumerate}
    \item Faça login no plano de locatário e encontre o grupo de
      usuários associado ao usuário afetado
    \item Adicione a permissão ECS FullAccess ao grupo de usuários
    \item Verifique se a permissão foi aplicada corretamente
  \end{enumerate}
  \end{workaroundsteps}

  \begin{verification}
  Faça login na página de locatário novamente, vá para a página do
  ECS, crie um ECS e verifique se a lista de projetos empresariais
  agora é exibida.
  \end{verification}

  \begin{rollback}
  Remova a permissão ECS FullAccess do grupo de usuários.
  \end{rollback}

  \begin{cleanup}
  Nenhuma limpeza necessária.
  \end{cleanup}

\end{workaround}

\section{Exemplos adicionais de formatação}

Esta seção demonstra comandos de formatação compartilhados disponíveis no modelo técnico.

\subsection{Caixas de destaque}

\begin{warning}
Esta é uma caixa de aviso. Use-a para destacar precauções importantes ou problemas potenciais.
\end{warning}

\begin{tip}
Esta é uma caixa de dica. Use-a para sugerir melhores práticas ou conselhos úteis.
\end{tip}

\begin{infobox}
Esta é uma caixa de informação. Use-a para informações gerais ou notas.
\end{infobox}

\subsection{Blocos de código}

\begin{code}[bash]
# Script de exemplo para verificação
kubectl get pods -n production
kubectl logs -f deployment/api-gateway -n production
\end{code}

\inlinecode{kubectl describe pod} para inspeção detalhada do pod.

\subsection{Tabelas}

\begin{hutable}{|l|c|c|}
\rowcolor{huaweired}
\thd{Componente} & \thd{Versão} & \thd{Status} \\
\tbody
API Gateway & 2.4.1 & OK \\
Banco de dados & 8.5.1 & OK \\
Cache & 7.0 & Degradado \\
\end{hutable}

\subsection{Comandos inline}

\note{Sempre verifique o backup antes de iniciar a solução alternativa.}

Referência: \weblink{https://docs.huaweicloud.com}{Documentação Huawei Cloud}

Navegação: \menu{Console > ECS > Instâncias}

Status: \badge{Verificado}

Parâmetro: \param{timeout} deve ser definido como 300 segundos.

\section{Paleta de Cores da Marca}

O template define cores auxiliares da marca alinhadas às
Diretrizes de Marca da Huawei Cloud:

\begin{hutable}{|l|l|l|l|}
\rowcolor{huaweired} \thd{Color} & \thd{HEX} & \thd{Pantone} & \thd{Sample} \\
\tbody
\textcolor{rosered}{Rose Red} & \texttt{\#C40054} & 7636 & \textcolor{rosered}{\rule{2cm}{4pt}} \\
\textcolor{darkred}{Dark Red} & \texttt{\#7F0001} & 483C & \textcolor{darkred}{\rule{2cm}{4pt}} \\
\textcolor{huaweiorange}{Orange} & \texttt{\#ED6D00} & 165C & \textcolor{huaweiorange}{\rule{2cm}{4pt}} \\
\textcolor{huaweiyellow}{Yellow} & \texttt{\#FCC800} & 7406C & \textcolor{huaweiyellow}{\rule{2cm}{4pt}} \\
\textcolor{huaweigreen}{Green} & \texttt{\#62B230} & 3501C & \textcolor{huaweigreen}{\rule{2cm}{4pt}} \\
\textcolor{huaweiblue}{Blue} & \texttt{\#30B5C5} & 2227C & \textcolor{huaweiblue}{\rule{2cm}{4pt}} \\
\end{hutable}

% --- Histórico de versões ---
\begin{changelog}
  \changelogentry{3.6.0}{\today}{
    \item Adicionado \inlinecode{\textbackslash setdocauthors} para demonstrar o recurso de autores na capa.
  }
  \changelogentry{1.3.0}{\today}{
    \item Atualização para modelo base v3.0.1 — correções menores (auto-descoberta de testes, textos de ajuda, comentários, CI).
  }
  \changelogentry{1.2.0}{\today}{
    \item Atualização para modelo base v3.0.0 — refatoração modular do pipeline (fábrica Lua compartilhada, lógica DOCX compartilhada, auto-descoberta do sistema de build).
  }
  \changelogentry{1.1.0}{\today}{
    \item Adicionada explicação em infobox e demonstrações de comandos de formatação compartilhados (caixas de destaque, código, tabelas, comandos inline).
  }
  \changelogentry{1.0.0}{\today}{
    \item Versão inicial.
  }
\end{changelog}

\end{document}
